\section{Introduction}

A guidance system makes two decisions, not one.

\begin{description}
\item[Level 1.] Which routing \emph{policy} should be active under the current
traffic conditions?
\item[Level 2.] Given that policy, which route does each guided vehicle
receive?
\end{description}

Level 2 is the subject of a large literature. This paper concerns Level 1, and
treats the policies themselves as fixed, published objectives.

The question matters because established policies embody different assumptions.
Distance minimisation assumes the shortest corridor can absorb what is sent to
it. Reactive travel-time routing assumes recent measurements predict the near
future, which fails when a large guided cohort acts on them
simultaneously~\cite{arnott1991information,acemoglu2018braess}. Load balancing
assumes capacity is unevenly distributed and worth equalising. Reliability-aware
routing assumes travel-time variability is itself a cost. Each assumption holds
in some conditions and not others, so the policy that should be active plausibly
depends on the traffic --- and production systems nevertheless fix it at design
time.

\paragraph{The gap this paper addresses.}
Showing that the preferred policy changes with conditions is not the same as
showing that a context-aware selector is worth building. A winner map can vary
substantially while the cost of ignoring it remains negligible, if the best
fixed policy loses only slightly wherever it is not best. These are separate
empirical quantities, and conflating them is easy because the first is the more
visible. We therefore measure both, and a third: whether a selector that
captures the available value can be trusted outside the conditions it was
trained on.

\paragraph{Research questions.}
\textbf{RQ1} Under which traffic, signal, information, penetration, capacity and
disruption conditions do established routing policies exhibit complementary
strengths?
\textbf{RQ2} Are the preference boundaries explicable by mechanistic traffic
variables rather than only by a statistical model?
\textbf{RQ3} Does a learned selector identify useful selection structure beyond
a strong mechanistic rule?
\textbf{RQ4} When the operating context lies outside the training support, or
uncertainty is too large, can the system abstain safely?

\paragraph{Contributions.}
\begin{enumerate}
\item A factorial characterisation of the operating regions in which four
established routing policies are complementary, over \NumContexts{} contexts
and \NumRunsMain{} paired microsimulation runs.
\item A mechanistic account of the preference boundary in dimensionless traffic
descriptors, reported as brackets between sampled levels.
\item A decision-oriented evaluation in which the primary comparison is a
learned selector against a mechanistic rule, not against a fixed policy, and in
which selection accuracy and decision regret are shown to rank methods
differently.
\item A selective mechanism with a support gate and a confidence gate,
evaluated separately under interpolation, extrapolation, concept shift and
domain shift, including a case it does not detect.
\end{enumerate}

We introduce no routing algorithm, no learning algorithm and no multi-criteria
decision method; every policy and model class used here is established and
cited. The contribution is the measurement and its interpretation. We make no
priority claim: each component below has precedent, and the positioning in
Section~\ref{sec:positioning} states what is and is not new.
